\section{Introduction}\label{sec:introduction}

A logarithmic Sobolev inequality controls the entropy of a function by
its Dirichlet energy.  Besides implying a Poincar\'e inequality, it gives
Gaussian concentration for Lipschitz functions.  For a log-concave
measure, a natural question is whether Gaussian tails of all
\emph{linear} functions already force this stronger conclusion.
We prove that they do, with a constant independent of dimension.

For a probability measure \(\nu\) and a nonnegative integrable function \(g\),
write
\[
 \Ent_\nu(g)=\int g\log g\,d\nu
       -\left(\int g\,d\nu\right)\log\left(\int g\,d\nu\right),
\]
with \(0\log0=0\).  A Lebesgue density \(\rho\) on \(\R^n\) is log-concave
if its positive set is convex and \(\log\rho\) is concave there.
We call a centered random vector \(X\) \emph{linearly \(a\)-subgaussian} if
\begin{equation}\label{eq:intro-subgaussian}
 \sup_{|\theta|=1}\E\exp\!\left(\frac{\langle X,\theta\rangle^2}{a^2}\right)
 \le2.
\end{equation}
This definition concerns only linear functions of \(X\).  We write \(D\)
for the Euclidean gradient.

\begin{theorem}\label{thm:main}
There is an absolute constant \(C<\infty\) with the following property.
For every \(n\ge1\), every centered log-concave probability measure
\(\mu\) with a Lebesgue density on \(\R^n\), and every \(a>0\)
satisfying \eqref{eq:intro-subgaussian} for \(X\sim\mu\),
\[
 \Ent_\mu(f^2)\le C a^2\int_{\R^n}|Df|^2\,d\mu
 \qquad\text{for every }f\in C_c^\infty(\R^n).
\]
The same constant works in every dimension.
\end{theorem}

No smoothness of the density or positive lower curvature bound is required.
In particular, the theorem includes uniform measures on convex bodies.
The proof will use smooth positive densities only after a reduction that
preserves any hypothetical sequence of counterexamples.

The Otto--Villani implication from logarithmic Sobolev to quadratic
transport--entropy inequalities \cite[Theorem~1]{OttoVillani2000}
gives the following standard consequence, with the same constant.
Write $\mathcal P_2(\R^n)$ for the Borel probabilities with finite
second moment. For $\alpha,\beta\in\mathcal P_2(\R^n)$, set
\[
 W_2(\alpha,\beta)^2
 =\inf_{\pi\in\Pi(\alpha,\beta)}\int|x-y|^2\,d\pi(x,y),
\]
where $\Pi(\alpha,\beta)$ is the set of their couplings. Relative
entropy is $H(\alpha\mid\beta)=\int f\log f\,d\beta$ if
$\alpha=f\beta$, with its possible value $+\infty$, and is
$+\infty$ if $\alpha\not\ll\beta$.

\begin{corollary}[Quadratic transport--entropy inequality]
\label{cor:intro-transport}
Under the assumptions of Theorem~\ref{thm:main}, every
$\nu\in\mathcal P_2(\R^n)$ satisfies
\[
 W_2(\nu,\mu)^2\le Ca^2 H(\nu\mid\mu),
\]
with the same absolute constant $C$ as in that theorem.
\end{corollary}

The proof in Section~\ref{sec:transport-consequence} gives the
approximation argument for nonsmooth log-concave densities.

\subsection{History and significance}

Gross's work on logarithmic Sobolev inequalities connected the
Gaussian inequality with hypercontractivity and established a
framework in which the constants remain independent of dimension
\cite{Gross1975}.  The diffusion approach of Bakry and \'Emery
\cite{BakryEmery1985} gives the inequality for a density $e^{-V}$
when $D^2V\ge\kappa\Id$ with $\kappa>0$: in the convention used
below, its constant is at most $\kappa^{-1}$.
Log-concavity alone asserts only $D^2V\ge0$ when the potential is
smooth.  The question addressed here is whether a global assumption
on the tails of linear functions can supply the logarithmic Sobolev
bound without positive curvature.

A logarithmic Sobolev inequality implies Gaussian concentration for
all Lipschitz functions through Herbst's exponential-moment argument;
see \cite[Section~2.3]{Ledoux1999}.  Under log-concavity, Gaussian
concentration for all Lipschitz functions also gives Gaussian
isoperimetry by Milman's theorem \cite[Theorem~1.1]{Milman2010},
and hence a logarithmic Sobolev inequality
\cite[Lemma~3.3]{Bobkov1999}, with dimension-independent bounds.
The question here asks whether control of linear functions suffices.
Uniformly subgaussian linear marginals are necessary, up to universal
constants, by the same Herbst argument.
Bizeul formulated their sufficiency within the log-concave class in 2023
\cite[Conjecture~2]{Bizeul2026}.  Theorem~\ref{thm:main} resolves
this conjecture for centered log-concave laws with a Lebesgue density,
under the precise marginal condition
\eqref{eq:intro-subgaussian}.

Bobkov's norm-based criterion
\cite[Theorem~1.3]{Bobkov1999} bounds the entropy coefficient by
a universal multiple of the squared subgaussian norm of $|X|$.
Under \eqref{eq:intro-subgaussian}, H\"older's inequality gives
$\E\exp(|X|^2/(na^2))\le2$, so that criterion yields a coefficient
of order $na^2$.  Bizeul's improvement
\cite[Theorem~4]{Bizeul2026} has dimension dependence $n^{1/4}$
in the parameter $\rho_{\mathrm{LS}}$ defined by
$\Ent(f^2)\le2\rho_{\mathrm{LS}}^2\int|Df|^2$.
Thus for linear subgaussian scale $a$ it gives an entropy coefficient
of order $\sqrt n\,a^2$.
Bizeul also obtained the dimension-free conclusion for rotationally
invariant log-concave laws \cite[Theorem~8]{Bizeul2026}.
Klartag and Lehec discuss the conjecture and these bounds in
\cite[Section~8.3, Conjecture~76 and Theorem~77]{KlartagLehec2024}.
The theorem's marginal hypothesis concerns more than covariance:
it controls the tails of every linear function, while its conclusion
controls the entropy of nonlinear tests.

Our preliminary calculus uses the Bakry--\'Emery semigroup argument,
including its entropy dissipation identity.  We give the anisotropic
and approximation forms needed for Gaussian posteriors in
Section~\ref{sec:calculus}.  The other classical input to that
calculus comes from Gaussian channels.  Guo, Shamai, and Verd\'u
relate the derivative of mutual information to minimum mean-square
estimation error \cite[Theorem~2]{GuoShamaiVerdu2005}; their
relative-entropy formulation \cite[Theorem~5, Equation~(66)]{GuoShamaiVerdu2005}
gives the Gaussian prediction identity used here.  We prove its
conditional integrated version with our variance and diffusion-time
conventions.  The entropy-loss comparison in
Lemma~\ref{lem:entropy-loss} requires an additional argument:
repeated observations of independent copies make the cost of discarding
conditional-mean directions negligible.  That comparison links
Gaussian prediction to the tensor near-equalities described next.

\subsection{Proof strategy}

Write \(R(\nu)\) for the best constant in the convention
\begin{equation}\label{eq:intro-lsi-convention}
 \Ent_\nu(f^2)\le2R(\nu)\int|Df|^2\,d\nu.
\end{equation}
Failure of Theorem~\ref{thm:main} would give, after truncation,
a weak quadratic tilt, and dilation, a sequence of compactly
supported log-concave laws \(\mu\) such that
\(R(\mu)=1\) while their linear subgaussian parameters tend to zero.
The argument derives incompatible descriptions of fluctuations
associated with these normalized laws.

\paragraph{From an extremal function to a predictable vector field.}
For a fixed \(s>0\), let \(G\) be a standard Gaussian independent of
\(X\sim\mu\).  An extremizer for the entropy inequality of
\(X+\sqrt sG\) yields a second probability law
\(\eta\) and a vector field
\[
 w=D\log\frac{d\eta}{d\mu_s},
 \qquad \mu_s=\Law(X+\sqrt s\,G).
\]
An exact change-of-reference identity shows that \(\eta\) itself
satisfies a logarithmic Sobolev inequality.
The Euler equation and translation tests make \(w\) an approximate
first eigenfunction of the diffusion associated with \(\eta\).
If the entropy constant instead coincides with the Poincar\'e
constant, take \(\eta=\mu_s\) and \(w=D\varphi\), where
\(\varphi\) is a gap eigenfunction.  Both cases are treated directly.

The crucial extra information concerns prediction of \(w(Y)\) from
a noisy observation of \(Y\sim\eta\).  With
\(I=\E_\eta|w|^2\) and a fresh independent standard Gaussian \(G\),
we prove the limiting identity
\[
 \frac{\E\big|\E[w(Y)\mid Y+\sqrt r\,G]\big|^2}{I}
 \longrightarrow e^{-r}
\]
on a fixed positive interval of noise variances.
It follows by matching entropy contraction with an entropy-loss
comparison proved using repeated Gaussian observations.
At each observation, we discard at most one direction: the difference
between the conditional means under two competing priors.
Independent copies make the accumulated cost of discarding these
directions negligible.

\paragraph{From prediction to tensor near-equalities.}
Gaussian polynomial orthogonality turns the exact prediction curve
into nearly sharp derivative estimates.
Starting with \(T_0=w\), we construct symmetric tensor fields \(T_l\)
whose derivatives approximate \(T_{l+1}\), all with asymptotically the
same squared norm and all concentrated near the first nonzero spectral value
of the diffusion.  The tensor estimates also control curvature and
exclude concentration in the relevant small-dimensional subspaces.

For each tensor \(T_l\), regard one index as a row and all remaining
indices as a column, and form the positive matrix \(T_lT_l^*\).
Averaging \(I^{-1}T_lT_l^*\) over a long band of levels gives a
positive matrix field \(\rho\) of asymptotic mean trace one.  Its square root is nearly an equality case in
Hilbert--Schmidt Poincar\'e.  Symmetry in the many tensor indices
forces the infinitesimal matrix coefficients to commute
asymptotically in a weighted Hilbert--Schmidt norm.  Gaussian linear combinations of these matrices, frozen at the initial
diffusion point, consequently have scalar Gaussian moments under the
weighted trace.  A Hermite expansion along the stationary diffusion
then gives covariance fluctuations with the lognormal factor
\[
 \exp(\sqrt{8t}\,Z-4t),\qquad Z\sim N(0,1),
\]
whose mean is one.

\paragraph{An information contradiction.}
Conditionally on the stationary point, observe a centered Gaussian
vector whose covariance is a bounded increasing function of
\(T\rho\), where \(T>0\) is a scale and the covariance cap is
adjustable.  Its marginal law is a mixture of these Gaussian laws.
Together with the square-root energy estimates, the logarithmic Sobolev
inequality of \(\eta\) gives a universal asymptotic upper bound on the
mutual information, integrated against \(dT/T^2\).
The matrix fluctuation law gives a lower bound by a one-dimensional
Gaussian mixture with the lognormal variance above.
After taking the sequence limit at fixed diffusion time and covariance
cap, that lower bound becomes arbitrarily large as these two
parameters increase together.

The scale integral is essential.  In a spectral block where
\(\rho\) has size \(\beta^2\), the substitution \(x=T\beta^2\) gives
\[
 \frac{dT}{T^2}=\beta^2\frac{dx}{x^2}.
\]
This is exactly the coefficient multiplying the unnormalized trace on
that block.
Consequently the information lower bound retains all of the matrix
mass even when the spectrum is arbitrarily diffuse.

Figure~\ref{fig:proof-dependencies} records the two constructions of the
critical law and the main dependencies of the information comparison.
\begin{figure}[tbp]
\centering
\begin{tikzpicture}[
  font=\small,
  stage/.style={draw, line width=.45pt, align=center,
    inner xsep=5pt, inner ysep=5pt},
  flow/.style={-{Stealth[length=2mm]}, line width=.55pt}
]
\node[stage, text width=6.4cm] (strict) at (-3.65,0)
  {Strict entropy branch\\
   Entropy optimizer\\
   Lemma~\ref{lem:critical-two-alternatives}(i)};
\node[stage, text width=6.4cm] (gap) at (3.65,0)
  {Gap branch\\
   Gap eigenfunction\\
   Lemma~\ref{lem:critical-two-alternatives}(ii)};

\node[stage, text width=5.8cm] (critical) at (0,-1.55)
  {Common critical law and vector field\\
   Proposition~\ref{prop:critical-sequence}};
\draw[flow] (strict.south) -- (critical.north west);
\draw[flow] (gap.south) -- (critical.north east);

\node[stage, text width=4.0cm, minimum height=1.5cm]
  (prediction) at (-4.85,-3.2)
  {Prediction curve\\
   and posterior estimates\\
   Prop.~\ref{prop:prediction}, Lemma~\ref{lem:posterior-outliers}};
\node[stage, text width=4.0cm, minimum height=1.5cm]
  (hierarchy) at (0,-3.2)
  {Symmetric\\tensor hierarchy\\
   Proposition~\ref{prop:hierarchy}};
\node[stage, text width=4.0cm, minimum height=1.5cm]
  (matrix) at (4.85,-3.2)
  {Matrix estimates\\and frozen law\\
   Proposition~\ref{prop:frozen-law}};
\draw[flow] (critical.west) -| (prediction.north);
\draw[flow] (prediction.east) -- (hierarchy.west);
\draw[flow] (hierarchy.east) -- (matrix.west);

\node[stage, text width=9.3cm] (contradiction) at (0,-5.45)
  {Integrated information contradiction\\
   Propositions~\ref{prop:information-upper}
   and~\ref{prop:information-lower},\\
   Lemma~\ref{lem:scalar-divergence}};
\draw[flow] (matrix.south) -- (4.85,-4.5)
  -- node[above, align=center] {matrix bounds} (0,-4.5)
  -- (contradiction.north);
\draw[flow] (critical.east) -- (7.45,-1.55)
  -- (7.45,-5.45)
  -- node[above, align=center] {entropy bound} (contradiction.east);
\end{tikzpicture}
\caption{Main stages of the contradiction argument. At the regularized
law, the strict entropy branch means that the logarithmic Sobolev
constant exceeds the Poincar\'e constant; in the gap branch the two
constants coincide. The logarithmic Sobolev inequality of the common
critical law, together with the square-root energy estimates of
Section~\ref{sec:matrices}, supplies the information upper bound.
The frozen matrix law supplies the scalar lower bound. After the
sequence limit at fixed parameters, its scale integral becomes
arbitrarily large as the diffusion time and covariance cap grow together.}
\label{fig:proof-dependencies}
\end{figure}


\subsection{Components of the argument}

Two ingredients have formulations beyond the final contradiction.
The entropy-loss comparison of Lemma~\ref{lem:entropy-loss} relates
relative entropy under common Gaussian convolution to conditional
prediction of the original relative score.  Its proof separates
the convexity input from the finite-copy observation argument.
The matrix reduction in Section~\ref{sec:matrices} converts symmetric
tensor near-equalities into a weighted scalar fluctuation law,
without requiring simultaneous diagonalization in varying
dimensions.  Section~\ref{sec:observations} then tests that law through
bounded covariance observations and integrates over scale.

For orientation, Section~\ref{sec:calculus} records the analytic
conventions and approximation facts.
Section~\ref{sec:critical} constructs the critical laws and vector
fields.
Section~\ref{sec:prediction} proves the prediction identity,
and Section~\ref{sec:hierarchy} constructs the tensor hierarchy.
Sections~\ref{sec:matrices} and \ref{sec:observations} turn that
hierarchy into the contradiction.

